% Part: first-order-logic
% Chapter: syntax-and-semantics
% Section: covered-structures

\documentclass[../../../include/open-logic-section]{subfiles}

\begin{document}

\olfileid{fol}{syn}{cov}

\olsection{પ્રથમ-ક્રમ ભાષાઓ માટેની આવૃત \printtoken{P}{structure}}

\begin{explain}
યાદ કરો કે કોઈ પદમાં એક પણ !!{variable}s ન હોય તો તે \emph{બંધ} છે.
\end{explain}

\begin{defn}[બંધ પદોનું !!^{value}]
જો $t$, ભાષા~$\Lang L$નું બંધ પદ હોય અને $\Struct M$,~$\Lang L$
માટેની !!{structure} હોય, તો \emph{!!{value}}~$\Value{t}{M}$ નીચે
પ્રમાણે વ્યાખ્યાયિત થાય છે:
\begin{enumerate}
\item જો $t$ માત્ર !!{constant}~$c$ હોય, તો $\Value{c}{M} = \Assign{c}{M}$.
\item જો $t$,~$\Atom{f}{t_1, \ldots, t_n}$ સ્વરૂપનું હોય, તો
  \[
  \Value{t}{M} = \Assign{f}{M}(\Value{t_1}{M}, \ldots,
  \Value{t_n}{M}).
  \]
\end{enumerate}
\end{defn}

\begin{defn}[આવૃત !!{structure}]
વ્યાપકક્ષેત્રનો દરેક ઘટક કોઈ બંધ પદનું !!{value} હોય તો !!{structure}
\emph{આવૃત} છે.
\end{defn}

\begin{ex}
!!{constant}s $\Obj{zero}$, $\Obj{one}$, $\Obj{two}$, \dots,
દ્વિસ્થાની !!{predicate}~$<$ અને દ્વિસ્થાની !!{function}s $+$ તથા
$\times$ ધરાવતી ભાષા~$\Lang L$ લો. ત્યારે~$\Lang L$ માટેની
!!a{structure}~$\Struct M$નું વ્યાપકક્ષેત્ર $\Domain M = \{0, 1, 2, \ldots
\}$ અને નિયુક્તિઓ $\Assign{\Obj{zero}}{M} = 0$,
$\Assign{\Obj{one}}{M} = 1$, $\Assign{\Obj{two}}{M} = 2$ વગેરે છે.
દ્વિસ્થાની સંબંધસંકેત~$<$ માટે,~$\Assign{<}{M}$ એ એવાં બધાં યુગ્મો
$\tuple{c_1, c_2} \in \Domain{M}^2$નો ગણ છે કે $c_1$,~$c_2$ કરતાં
નાનો હોય: દાખલા તરીકે, $\tuple{1, 3} \in \Assign{<}{M}$ પરંતુ
$\tuple{2, 2} \notin \Assign{<}{M}$. દ્વિસ્થાની !!{function}~$+$
માટે $\Assign{+}{M}$ને સામાન્ય રીતે વ્યાખ્યાયિત કરો---દાખલા તરીકે,
$\Assign{+}{M}(2,3)$નું મૂલ્ય~$5$ થાય છે---અને દ્વિસ્થાની
!!{function}~$\times$ માટે પણ એવું જ કરો. તેથી $\Obj{four}$નું
!!{value} માત્ર~$4$ છે અને $\times(\Obj{two},
+(\Obj{three},\Obj{zero}))$નું !!{value} (અથવા મધ્યસ્થ સંજ્ઞામાં,
$\Obj{two} \times (\Obj{three} + \Obj{zero})$) આ છે:
\begin{multline*}
\Value{\times(\Obj{two}, +(\Obj{three},\Obj{zero}))}{M} =\\
\begin{aligned}
& =\Assign{\times}{M}(\Value{\Obj{two}}{M}, \Value{+(\Obj{three}, \Obj{zero})}{M})\\
& = \Assign{\times}{M}(\Value{\Obj{two}}{M}, \Assign{+}{M}(\Value{\Obj{three}}{M},
\Value{\Obj{zero}}{M})) \\
& = \Assign{\times}{M}(\Assign{\Obj{two}}{M}, \Assign{+}{M}(\Assign{\Obj{three}}{M},
\Assign{\Obj{zero}}{M})) \\
& = \Assign{\times}{M}(2, \Assign{+}{M}(3, 0)) \\
& = \Assign{\times}{M}(2, 3) \\
& = 6
\end{aligned}
\end{multline*}
\end{ex}

\begin{prob}
અંકગણિતનો પ્રમાણભૂત નિદર્શ~$\Struct N$ આવૃત છે? સમજાવો.
\end{prob}

\end{document}
